\documentclass[a4paper,12pt]{article}

% ======================
% Encodage / langue
% ======================
\usepackage[utf8]{inputenc}
\usepackage[T1]{fontenc}
\usepackage[french]{babel}

% ======================
% Mise en page / style
% ======================
\usepackage[margin=1in]{geometry}
\usepackage{graphicx}
\usepackage{xcolor}
\usepackage{tikz}
\usepackage{titlesec}
\usepackage{tocloft}
\usepackage{setspace}
\usepackage{float}
\usepackage{placeins}
\usepackage{xurl}
\usepackage{hyperref}
\usepackage{enumitem}
\usepackage{booktabs}

% ======================
% Couleurs
% ======================
\definecolor{ensiared}{RGB}{156,31,46}
\definecolor{tocblue}{RGB}{0,102,204}

% ======================
% ToC en bleu
% ======================
\renewcommand{\cftsecfont}{\color{tocblue}}
\renewcommand{\cftsubsecfont}{\color{tocblue}}
\renewcommand{\cftsecpagefont}{\color{tocblue}}
\renewcommand{\cftsubsecpagefont}{\color{tocblue}}

% ======================
% Interligne / spacing
% ======================
\linespread{1.15}
\titlespacing*{\section}{0pt}{*2}{*1.2}
\titlespacing*{\subsection}{0pt}{*1.5}{*1}

% ======================
% Macros
% ======================
\newcommand{\todo}{\textit{(À compléter)}}

% Encadré robuste (évite les erreurs tabularx)
\usepackage[most]{tcolorbox}
\newtcolorbox{tipbox}{
  colback=black!3,
  colframe=black!25,
  boxrule=0.6pt,
  arc=2pt,
  left=6pt,right=6pt,top=6pt,bottom=6pt
}

% Capture (optionnelle)
\newcommand{\screenshot}[2][]{%
  \begin{figure}[H]
    \centering
    \fbox{\includegraphics[width=0.95\textwidth]{#2}}
    \caption{#1}
  \end{figure}
}

\begin{document}

% ==========================================================
% Page de garde
% ==========================================================
\pagenumbering{gobble}

\begin{titlepage}
  \begin{tikzpicture}[remember picture, overlay]
    \fill[ensiared] ([yshift=2cm]current page.south west) rectangle (current page.south east);
    \fill[ensiared] ([yshift=2cm]current page.north west) rectangle (current page.north east);
  \end{tikzpicture}

  \centering
  \vspace*{2cm}

  % Logos
  \includegraphics[width=0.35\textwidth]{img/sd logo.jpg}\hfill
  \includegraphics[width=0.25\textwidth]{img/uca logo.png}\hfill

  \vspace*{1.0cm}

  {\color{white}\scshape\LARGE Base de données \& Systèmes décisionnels\par}
  \vspace{0.4cm}
  {\scshape\Large UNIVERSITÉ CLERMONT AUVERGNE — Campus Aurillac\par}

  \vspace{1.2cm}
  \hrule height 1pt
  \vspace{0.5cm}
  {\huge\bfseries Guide utilisateur\par}
  \vspace{0.2cm}
  {\Large Outil décisionnel — \textit{Ventes \& Logistique} (Power BI)\par}
  \vspace{0.5cm}
  \hrule height 1pt
  \vspace{1.6cm}

  \begin{minipage}[t]{0.5\textwidth}
    \raggedright
    {\Large\textbf{Réalisé par} :\par}
    \vspace{0.2cm}
    \begin{tabular}{l}
      \ Ibrahima BODIAN \\
    \end{tabular}
  \end{minipage}%
  \begin{minipage}[t]{0.5\textwidth}
    \raggedleft
    {\Large\textbf{Encadrant} :\par}
    \vspace{0.2cm}
    \textit{\ Bruno LAGARDE} \\
  \end{minipage}

  \vfill
  {\large Date : \textit{\today}\par}
\end{titlepage}

% ==========================================================
% Sommaire
% ==========================================================
\renewcommand{\contentsname}{Sommaire}
\tableofcontents
\newpage
\pagenumbering{arabic}
\setcounter{page}{1}

% ==========================================================
% 1) Objectif
% ==========================================================
\section{Objectif de l’outil}
Ce rapport Power BI permet d’analyser :
\begin{itemize}[leftmargin=*]
  \item le chiffre d’affaires (CA), les quantités vendues, les remises et les frais de port ;
  \item des vues de pilotage par produits/catégories, clients, zones (pays/ville) et transporteurs ;
  \item la logistique (écarts date objectif vs date d’envoi, réapprovisionnement, quantités livrées par zone).
\end{itemize}

\begin{tipbox}
\textbf{À retenir :} ce guide décrit \textbf{l’utilisation du rapport Power BI} (navigation, filtres, lecture des visuels, export).
Les détails techniques (DW, scripts, contrôles SQL) sont traités dans la \textbf{Fiche technique}.
\end{tipbox}

% ==========================================================
% 2) Navigation + F/L
% ==========================================================
\section{Navigation dans le rapport (pages et références F/L)}

\subsection{Pages du rapport}
\begin{itemize}[leftmargin=*]
  \item \textbf{Page 1 — Synthèse Finance} : indicateurs clés et vue globale.
  \item \textbf{Page 2 — Produits \& Catégories (F2, F6)}
  \item \textbf{Page 3 — Clients \& Remises (F1, F5, F8)}
  \item \textbf{Page 4 — Transport \& Coûts (F3, F4, F7)}
  \item \textbf{Page 5 — Logistique (L1–L5)}
\end{itemize}

\subsection{Correspondance avec la consigne (Finance F)}
\begin{table}[H]
\centering
\renewcommand{\arraystretch}{1.2}
\begin{tabular}{p{2cm} p{11.5cm}}
\toprule
\textbf{Réf.} & \textbf{Où la consulter dans Power BI} \\
\midrule
F1 & Page 3 : distribution des remises + détail par client/produit/employé \\
F2 & Page 2 : CA par produit + lecture par catégorie \\
F3 & Page 4 : CA par pays de livraison + comparaison avec frais de port \\
F4 & Page 4 : CA / expéditions par transporteur \\
F5 & Page 3 : commandes par employé + remise moyenne par client \\
F6 & Page 2 : CA et quantités par fournisseur et produit \\
F7 & Page 4 : port total + nombre d’expéditions par transporteur \\
F8 & Page 3 : top clients / concentration du CA (rapport ajouté) \\
\bottomrule
\end{tabular}
\caption{Correspondance Finance (F) et pages du rapport.}
\end{table}

\subsection{Correspondance avec la consigne (Logistique L)}
\begin{table}[H]
\centering
\renewcommand{\arraystretch}{1.2}
\begin{tabular}{p{2cm} p{11.5cm}}
\toprule
\textbf{Réf.} & \textbf{Où la consulter dans Power BI (Page 5)} \\
\midrule
L1 & Table \textit{Plus grands écarts Date objectif vs Date envoi} (tri sur l’écart) \\
L2 & Table \textit{Produits à réapprovisionner} (stock dispo vs niveau réap.) \\
L3 & Visuel \textit{Quantités livrées par région et ville} \\
L4 & (Si présent) liste liée aux commandes/produits/fournisseurs selon votre implémentation \\
L5 & Indicateur/rapport utile au décideur logistique (ex. délai moyen, nb expéditions, etc.) \\
\bottomrule
\end{tabular}
\caption{Correspondance Logistique (L) et visuels de la Page 5.}
\end{table}

% ==========================================================
% 3) Filtres et sélections
% ==========================================================
\section{Filtres et sélections (slicers)}

\subsection{Utiliser les segments (slicers)}
\begin{itemize}[leftmargin=*]
  \item Cliquez sur une valeur (ex. \textit{Pays livraison}, \textit{Semaine}) pour filtrer la page.
  \item Pour sélectionner plusieurs valeurs : maintenez \textbf{Ctrl} pendant les clics.
  \item Pour effacer rapidement une sélection : utilisez l’icône \textbf{gomme} du segment.
\end{itemize}

\subsection{Bon réflexe si ``tout semble vide''}
\begin{itemize}[leftmargin=*]
  \item Vérifiez si un segment est trop restrictif (pays, semaine, transporteur).
  \item Effacez la sélection (gomme) pour revenir à la vue globale.
\end{itemize}

% ==========================================================
% 4) Comprendre et explorer les graphiques
% ==========================================================
\section{Comprendre et explorer les graphiques}
\subsection{Lecture du nuage de points \textit{Prix vs Quantité} (Page 2)}
\begin{itemize}[leftmargin=*]
  \item À droite : produits plus chers ; en haut : produits vendus en plus grande quantité ; bulles plus grosses : plus de CA.
  \item Objectif : distinguer produits ``volume'' vs produits ``premium''.
\end{itemize}

\subsection{Lecture des tables logistiques (Page 5)}
\begin{itemize}[leftmargin=*]
  \item \textbf{Écart objectif/envoi} :
  \begin{itemize}
    \item écart \textbf{négatif} = envoi \textbf{avant} la date objectif (avance) ;
    \item écart \textbf{positif} = envoi \textbf{après} la date objectif (retard).
  \end{itemize}
  \item \textbf{Réapprovisionnement} : un produit est à réapprovisionner lorsque \textit{stockDispo} est inférieur au \textit{niveau réap.}.
\end{itemize}

% ==========================================================
% 5) Mode d’emploi page par page
% ==========================================================
\section{Mode d’emploi par page}

\subsection*{Page 1 — Synthèse Finance}
\begin{itemize}[leftmargin=*]
  \item Objectif : obtenir une vue rapide (CA, commandes, quantités).
\end{itemize}

\subsection*{Page 2 — Produits \& Catégories (F2, F6)}
\begin{itemize}[leftmargin=*]
  \item \textbf{F2} : identifier les produits/catégories qui contribuent le plus au CA.
  \item \textbf{F6} : analyser la contribution des fournisseurs (matrix fournisseur $\rightarrow$ produit).
\end{itemize}

\subsection*{Page 3 — Clients \& Remises (F1, F5, F8)}
\begin{itemize}[leftmargin=*]
  \item \textbf{F8} : repérer les clients les plus contributeurs (top clients).
  \item \textbf{F1} : comprendre la distribution et l’impact des remises.
  \item \textbf{F5} : comparer les commandes/indicateurs par employé et par client.
\end{itemize}

\subsection*{Page 4 — Transport \& Coûts (F3, F4, F7)}
\begin{itemize}[leftmargin=*]
  \item \textbf{F3} : comparer les pays de livraison (CA/port).
  \item \textbf{F4/F7} : comparer les transporteurs (activité, coûts, etc.).
\end{itemize}

\subsection*{Page 5 — Logistique (L1–L5)}
\begin{itemize}[leftmargin=*]
  \item \textbf{Slicers principaux} : \textit{Pays livraison} et \textit{Semaine}.
  \item \textbf{KPI} : \textit{Nb expéditions} = nombre de commandes/expéditions sur le périmètre filtré.
  \item \textbf{L1} : table des plus grands écarts (tri sur l’écart absolu).
  \item \textbf{L2} : table des produits à réapprovisionner (stock dispo, seuil).
  \item \textbf{L3} : quantités livrées par région et ville.
  \item \textbf{L5} : indicateurs utiles (ex. volume d’expéditions / délais / écarts) pour orienter la décision.
\end{itemize}

% ==========================================================
% 6) Exporter
% ==========================================================
\section{Exporter et partager}

\subsection{Exporter le rapport en PDF (Power BI Desktop)}
\begin{itemize}[leftmargin=*]
  \item \textbf{Fichier} $\rightarrow$ \textbf{Exporter} $\rightarrow$ \textbf{Exporter au format PDF}.
\end{itemize}


% ==========================================================
% Annexes
% ==========================================================
\newpage
\appendix

\section{Glossaire}
\begin{itemize}[leftmargin=*]
  \item \textbf{CA net} : CA après remise.
  \item \textbf{Remise} : réduction appliquée (ex. 0.05 = 5\%).
  \item \textbf{Port} : frais de transport au niveau de la commande.
  \item \textbf{Nb expéditions} : nombre de commandes/expéditions (grain commande).
  \item \textbf{Stock dispo} : unités en stock moins unités en commande.
  \item \textbf{Niveau réap.} : seuil à partir duquel un réapprovisionnement est recommandé.
  \item \textbf{Écart objectif/envoi} : différence en jours entre date objectif et date d’envoi.
\end{itemize}

\section{FAQ (dépannage rapide)}
\begin{itemize}[leftmargin=*]
  \item \textbf{Je ne vois plus de données} : vérifier les slicers, puis effacer la sélection (gomme).
  \item \textbf{Je veux vérifier un chiffre} : exporter le visuel en CSV et contrôler les lignes filtrées.
  \item \textbf{Trop d’informations sur une table} : appliquer un tri ou un filtre Top N (ex. Top 10).
\end{itemize}

\end{document}